% ============================================================
% §4.4  Factuality Assessment
% ============================================================

\subsection{Factuality Assessment}
\label{sec:factuality}

Factual accuracy is measured via FActScore~\cite{min2023factscore}, which
decomposes each generated response into atomic claims and verifies each claim
against the domain-partitioned knowledge base used during retrieval
(Section~\ref{ssec:setup}), yielding a \emph{hallucination rate} defined
as the fraction of unverified claims per response.%
\footnote{%
  We initially conducted G-Eval~\cite{liu2023geval} quality assessment using
  Qwen2.5-72B-Instruct as an independent judge over 50 sampled CROSS-domain
  queries; results are retained in Appendix~\ref{appendix:geval} as
  preliminary reference.
  Subsequent attempts to extend this evaluation using
  Llama-3.1-70B-Instruct ($n{=}200$) revealed a systematic ceiling effect
  (92.5\% of scores concentrated at 9/10 on a 1--10 scale,
  inter-model range $<$0.3 points), indicating insufficient discriminative
  resolution for fine-grained model comparison.
  We therefore adopt FActScore-based hallucination assessment as the primary
  factuality metric.
}

% --- Baseline results ---
Table~\ref{tab:factscore_baselines} presents hallucination rates for all five
7B-backbone models.
NexRoute achieves the lowest overall hallucination rate of 8.9\%---a 16.0\%
relative reduction compared to Naive RAG (10.6\%).
This reduction is statistically significant relative to Adaptive RAG
($p{=}0.001$) and CRAG ($p{=}0.026$), and represents a consistent trend
relative to Naive RAG ($p{=}0.075$), though the latter does not reach
conventional significance at $\alpha{=}0.05$.
Domain-level breakdown reveals that NexRoute's advantage spans both evaluation
subsets: the CROSS-domain hallucination rate is 4.3\%---tied with Naive RAG
as the lowest among all models---while the single-domain rate is 11.7\%,
below all other 7B baselines (Naive RAG: 12.6\%).
Adaptive RAG and CRAG exhibit overall hallucination rates above Naive RAG
(12.1\% and 11.1\%, respectively), suggesting that their adaptive retrieval
heuristics introduce additional factual noise without corresponding accuracy
gains in this multi-domain setting.

% --- Ablation: V02 ---
The ablation study (Table~\ref{tab:factscore_ablation}) identifies the
\textbf{Function Axis} as the dominant structural contributor to factuality.
Removing the Function Axis (V02) increases overall hallucination from 8.9\%
to 15.2\% ($+$6.3\,pp; $p{<}0.001$), with significant degradation in both
CROSS ($+$4.8\,pp; $p{=}0.003$) and single-domain ($+$6.7\,pp; $p{<}0.001$)
subsets.
The Function Axis partitions inference responsibility among LoRA-based domain
experts via function-typed routing; its removal forces all queries through a
single generalist pathway, substantially increasing the rate of unverified
claims.

% --- Ablation: V03 ---
Removing Partitioned Retrieval (V03), which implements the Isolated Domain
Index (IDI) mechanism, produces a significant increase in single-domain
hallucination from 11.7\% to 16.7\% ($+$5.0\,pp; $p{=}0.003$), confirming
that the IDI is the primary structural driver of factuality in specialized
domain contexts.
The effect on CROSS-domain queries ($+$1.8\,pp; $p{=}0.19$) does not reach
statistical significance, consistent with the observation that CROSS queries
draw evidence from multiple domain indices simultaneously, distributing---and
thereby diluting---the contribution of any single partitioned index.

% --- V09 ---
Among CROSS-domain merging strategies, the Dominant-Single approach (V09)
yields a CROSS hallucination rate of 9.9\%---more than doubling V01---
indicating that collapsing cross-domain evidence onto a single-domain expert
substantially degrades factual precision on interdisciplinary queries.

